\documentclass[10pt,a4paper]{amsart}
\usepackage[margin=19mm]{geometry}
\usepackage{amsmath,amssymb,mathtools}
\usepackage[scr=rsfs,cal=euler]{mathalfa}
\usepackage{xfrac}
\usepackage[unicode,hidelinks,bookmarks=false]{hyperref}
\newcommand{\ie}{i.e.\ }
\newcommand{\cf}{cf.\ }
\newcommand{\resp}{respectively\ }
\newcommand{\Subsection}{\subsection}
\providecommand{\nobreakdash}{\nobreak\hbox{-}}
\setlength{\emergencystretch}{2em}
\multlinegap0pt
\usepackage{bm}
\usepackage{bbm}
\usepackage{dsfont}
\usepackage{xspace}
\usepackage{cancel}
\usepackage{mathtools}
\usepackage{amsthm}
\makeatletter
\@ifundefined{theorem}{\newtheorem{theorem}{Theorem}}{}
\@ifundefined{corollary}{\newtheorem{corollary}[theorem]{Corollary}}{}
\@ifundefined{lemma}{\newtheorem{lemma}[theorem]{Lemma}}{}
\@ifundefined{proposition}{\newtheorem{proposition}[theorem]{Proposition}}{}
\makeatother
\title[Quasi-maximal ideals and ring extensions]{Quasi-maximal ideals and ring extensions}
\author{Gabriel Picavet, Martine Picavet-L'Hermitte}
\date{Source: arXiv:2601.13093v1 (CC BY 4.0)}
\begin{document}
\maketitle

\noindent\emph{Source and licence.} Quasi-maximal ideals and ring extensions. Version arXiv:2601.13093v1 (CC BY 4.0), distributed under the Creative Commons Attribution 4.0 International licence, as recorded in the arXiv licence metadata for this exact version and shown on the abstract page https://arxiv.org/abs/2601.13093v1. Full rights notice: licenses/arxiv-2601.13093-CC-BY-4.0.txt.\par\smallskip

\noindent\textit{Editorial scope.} Counted unit: the Proposition whose source label is \texttt{3.52} in the article (source0.tex lines 407--417, source label \texttt{3.52}), delivered together with the article's own complete original proof of it (source0.tex lines 418--435, one \texttt{proof} environment of 18 lines). No mathematics is written, repaired or re-derived: statement and proof are the article's own text line for line. Required context retained uncounted: 3.2 (lines 325--326); 3.53 (lines 394--399). Every definition, display and label of those lines is retained unchanged. Omitted from this package and not redistributed: 1--324, 327--393, 400--406, 436--1271. Nothing inside a retained excerpt is omitted or altered: every retained body is delivered exactly as the article's lines are written, so no omission receipt of any kind accompanies this package. Citations: the retained excerpts carry no citation command at all, so no bibliography entry of the article is transcribed and no reference is redistributed. Reference adaptation: none was needed and none was made; every reference inside the delivered unit resolves inside it, so no unresolved reference or citation is produced. Printed slips kept literal: no printed word, symbol, hypothesis or number of the counted statement or of its proof was altered. Layout and class: the article's own class is \texttt{amsart} with packages including amsmath, amssymb, amscd, amsfonts, stmaryrd, turnstile, mathrsfs, eucal, color, hyperref; the wrapper is the lane's pinned \texttt{amsart} editorial wrapper at 10pt on A4, which restates the theorem-like environments the retained text needs and adds the article's own macro definitions that the retained text needs (none), with extra wrapper packages (bm, bbm, dsfont, xspace, cancel, mathtools). The delivered numbering is the unit's own. Assets: no retained span contains a figure environment, a graphics-inclusion command, a PDF-page inclusion, a table or a TikZ picture; no image file is copied into this package, the package declares no asset dependency and no image file is redistributed with it. Licence: version arXiv:2601.13093v1 of this article is distributed under the Creative Commons Attribution 4.0 International licence, as recorded in the arXiv licence metadata for this exact version and shown on the abstract page https://arxiv.org/abs/2601.13093v1. A full rights notice is delivered at licenses/arxiv-2601.13093-CC-BY-4.0.txt. Characters: every retained excerpt is pure ASCII, so the delivered engine, XeLaTeX with its default Latin Modern text font, reports no missing character.
\begin{corollary}\label{3.2} Let $R\subseteq S$ be an extension and $I\in \mathrm{Q_rMax}(R)$ be such that $\sqrt[R]IS\in\mathrm{Max}(S)$. Then $I=IS\cap R$ if and only if $IS\in\mathrm {Q_rMax}(S)$. 
\end{corollary}

\begin{lemma}\label{3.53} Let $f:R\to S$ be a ring morphism, $J$ an ideal of $S$ and $I:=f^{-1}(J)$. Then, $I=f^{-1}(I\cdot S)$.
 \begin{enumerate}
\item There exists an ideal $K$ of $S$ maximal in $\mathcal F:=\{K'$ ideal of $S\mid J\subseteq K',\ I=f^{-1}(K')\}$. Such an ideal $K$ is called a Max-upper of $I$ containing $J$.
\item If $N\in\mathrm{Max}(S)\cap\mathrm{V}_S(J)$, there exists an ideal $K_1$ of $S$ maximal in $\mathcal F':=\{K'$ ideal of $S\mid J\subseteq K' \subseteq N,\ I=f^{-1}(K')\}$. Such an ideal $K_1$ is called a Max-upper of $I$ containing $J$ and contained in $N$.
\end{enumerate}
 \end{lemma}

\begin{proposition}\label{3.52} Let $R\subseteq S$ be an extension, $J$ an ideal of $S$ and $I:=J\cap R$. Assume that $I\prec M$ for some $M\in\mathrm{Max}(R)$ and $N:=MS\in\mathrm{Max}(S)$ with $J\subseteq N$. Then, there exists a Max-upper $K$ of $S$ of $I$ containing $J$ and contained in $N$. Moreover, $K\in\mathrm{QMax}(S)$ and the next 
  properties hold:
\begin{enumerate} 
\item If $I\in\mathrm{Q_rMax}(R)$, then $K=IS\in\mathrm{Q_rMax}(S)$. 
\item Let $I\in\mathrm{Q_dMax}(R)$ with $I=M_1\cap M_2,\ M_i\in\mathrm {Max}(R)$ for $i\in\mathbb N_2$. Set $N_i:=M_iS$ for $i\in\mathbb N_2$, and assume that $N_1\in\mathrm{Max}(S)$. Then $N_2\neq S$ and we have two cases: 
\begin{enumerate}
\item If $N_2\in\mathrm{Max}(S)$, then $K=IS\in\mathrm{Q_dMax}(S)$.  
\item If $N_2\not\in\mathrm{Max}(S)$, then $IS\subset K$ and $IS\not\in \mathrm{QMax}(S)$.
\end{enumerate}
\end{enumerate}
\end{proposition}

\begin{proof}   
Assume that $I\prec M$ for some $M\in\mathrm{Max}(R)$ such that $N:=M S\in\mathrm{Max}(S)$ with $J\subseteq N$. Then, $I\in\mathrm{Q_rMax}(R)\cup\mathrm{Q_dMax}(R)$. According to Lemma \ref{3.53}, we get that $I=IS\cap R$ and there exists an ideal $K$ of $S$ Max-upper of $I$ containing $J$ and contained in $N$. We claim that $K\in\mathrm{QMax}(S)$. First, $K\subset N$ since $I\prec M$. Assume that $K\nprec N$. Hence, there exists some ideal $K''$ of $S$ such that $K\subset K''\subset N$, so that $I\subseteq K''\cap R\subseteq M=N\cap R$. But $I\prec M$ gives that either $I=K''\cap R\ (*)$ or $K''\cap R=M\ (**)$. In case $(*)$, we have $K''\in\mathcal F$, a contradiction with the maximality of $K$. In case $(**)$, we have $M\subseteq K''$, giving $N=MS\subseteq K''$, also a contradiction. To conclude, $K\prec N$, whence $K\in\mathrm{QMax}(S)$.

(1) If $I\in\mathrm{Q_rMax}(R)$, $M$ is the only maximal ideal of $R$ with $I\subset M$ and $MS\in\mathrm{Max}(S)$. As $IS\subseteq K\subset N= MS$, we have $IS\neq\sqrt[S]{IS}$. According to Corollary \ref{3.2}, 
 $IS\in\mathrm{Q_rMax}(S)$ and $K=IS$. 

(2) Assume that $I\in\mathrm{Q_dMax}(R)$ and set $I=M_1\cap M_2,\ M_i\in\mathrm{Max}(R)$ for $i\in\mathbb N_2$, with $N_i:=M_iS$. Set, for instance, $M=M_1$, so that $N_1=N\in\mathrm{Max}(S)$. Then, $IS=M_1 M_2S=NN_2$. As $M_1$ and $M_2$ are comaximal in $R$, so are $N$ and $N_2$ in $S$. 

First, if $N_2=S$, then $I=M_1\cap M_2=M_1M_2$ implies $IS =M_1S=N$ which shows that $I=IS\cap R=N\cap R=M$, a contradiction: this case cannot appear. We  consider the two possibilities for $N_2$.

(a) If $N_2\in\mathrm{Max}(S)$, the same reasoning as for $N$ shows that $IS=N_1\cap N_2$ and $IS=K\in\mathrm {Q_dMax}(S)$. 

(b) Assume that $N_2\not\in\mathrm{Max}(S)$. It follows that $IS= N\cap N_2$. As $K\prec N$, either $K\in\mathrm{Q_dMax}(S)$ ($\alpha$) or $K\in\mathrm {Q_rMax}(S)$ ($\beta$). 

In case ($\alpha$), $K=N\cap N'$, where $N'\neq N$ and $N'\in\mathrm {Max}(S)$ with $IS\subseteq N_2\subset N'$. In this case, $IS\subset K$ and $IS$ is not in $\mathrm{QMax}(S)$.

In case ($\beta$), $N$ is the only maximal ideal of $S$ with $K\subseteq N$. But $IS=NN_2$ shows that $IS$ is contained in at least two maximal ideals of $S$. Then, $IS\subset K$ and $IS$ is not in $\mathrm{QMax}(S)$. 
\end{proof}

\end{document}
